\section{Baccalauréat séries C/E — Session 2022}
\noindent\begin{minipage}{0.5\linewidth}\footnotesize Office du Baccalauréat du Cameroun\end{minipage}%
\hfill\begin{minipage}{0.45\linewidth}\footnotesize\raggedleft Examen : \textbf{Baccalauréat} · Série : \textbf{C/E}\\
Épreuve : \textbf{Mathématiques} · Durée : \textbf{4\,h} · Coef. : \textbf{7 -- 6}\end{minipage}
\par\smallskip{\color{onbo}\rule{\linewidth}{1pt}}

\noindent\textit{Cette épreuve est constituée de deux parties indépendantes.}

\subsection*{Partie A — Évaluation des ressources \hfill (15 points)}

\noindent\textbf{Exercice 1.} \hfill\textit{(5 points)}\\
Le plan complexe est muni d'un repère orthonormé direct $(O\,;\,\vec{e_1},\vec{e_2})$. On
considère les points $A$, $B$, $F$ et $G$ d'affixes respectives :
\[ Z_A=1+\mathrm{i}\sqrt3\ ;\quad Z_B=-1-\mathrm{i}\sqrt3\ ;\quad Z_F=4\ ;\quad Z_G=-4. \]
\begin{enumerate}[label=\arabic*)]
\item Résoudre dans $\C$ l'équation $z^2+2-2\mathrm{i}\sqrt3=0$.
\item Soit $s$ la similitude directe d'expression complexe $z'=(1-\mathrm{i}\sqrt3)\,z$.
  \begin{enumerate}[label=\textbf{\alph*.}]
  \item Donner les éléments caractéristiques de $s$.
  \item Quelles sont les images par $s$ des points $A$ et $B$ ?
  \end{enumerate}
\item Soit $(\mathcal{E})$ l'ellipse de foyers $A$ et $B$ et d'excentricité $e=\frac12$.
  \begin{enumerate}[label=\textbf{\alph*.}]
  \item Déterminer une équation de l'image $(\mathcal{E}')$ de $(\mathcal{E})$ par la similitude $s$.
  \item Construire $(\mathcal{E}')$ puis $(\mathcal{E})$ dans le même repère.
  \end{enumerate}
\item Aïcha a choisi au hasard l'un après l'autre, deux points distincts parmi les points
  $O$, $A$, $B$, $F$ et $G$ comme ceux par lesquels passe l'axe focal de l'ellipse $(\mathcal{E}')$.
  Quelle est la probabilité qu'elle ait choisi deux points de l'axe focal de $(\mathcal{E}')$ ?
\end{enumerate}

\smallskip
\noindent\textbf{Exercice 2.} \hfill\textit{(5 points)}\\
Soit $(\vec\imath,\vec\jmath,\vec k)$ une base d'un espace vectoriel $E$. Soit $f$ un endomorphisme de $E$.
\begin{enumerate}[label=\arabic*)]
\item Pour $k$ appartenant à $\R$, on considère l'ensemble $E_k$ des vecteurs $\vec u$ de $E$
  tels que $f(\vec u)=k\vec u$.
  \begin{enumerate}[label=\textbf{\alph*.}]
  \item Démontrer que $E_k$ est un sous-espace vectoriel de $E$.
  \item On suppose que $f$ vérifie l'égalité $f\circ f=2f$. Démontrer que $\vec u\in\operatorname{Im}f$
    si et seulement si $\vec u\in E_2$.
  \end{enumerate}
\item On suppose ici qu'on a :
  \[
  \left\{\begin{aligned}
  f(\vec\imath+\vec\jmath)&=2\vec\imath+2\vec\jmath\\
  f(\vec\imath-\vec\jmath)&=2\vec\imath-2\vec\jmath\\
  f(\vec\imath-\vec\jmath+\vec k)&=\vec 0.
  \end{aligned}\right.
  \]
  \begin{enumerate}[label=\textbf{\alph*.}]
  \item Démontrer que $f(\vec\imath)=2\vec\imath$, $f(\vec\jmath)=2\vec\jmath$ et $f(\vec k)=-2\vec\imath+2\vec\jmath$.
  \item Donner la matrice $M$ de $f$ dans la base $(\vec\imath,\vec\jmath,\vec k)$.
  \item Démontrer que $f\circ f=2f$.
  \item Déterminer, par une de ses bases, le noyau $\operatorname{Ker}f$ de $f$.
  \item Déterminer l'image $\operatorname{Im}f$ de $f$. On précisera une de ses bases.
  \end{enumerate}
\end{enumerate}

\smallskip
\noindent\textbf{Exercice 3.} \hfill\textit{(5 points)}\\
$f$ est une fonction définie sur $[0\,;\,2\pi]$ par $f(x)=\mathrm{e}^{-x}\cos(x)$. $(\mathcal{C}_f)$ est
la courbe de $f$ dans un repère orthogonal où en abscisse on a $2$~cm pour unité et en
ordonnée $4$~cm pour unité.
\begin{enumerate}[label=\arabic*)]
\item Démontrer que $f''(x)+2f'(x)+2f(x)=0$.
\item Étudier les variations de $f$ et dresser son tableau de variations.
\item \begin{enumerate}[label=\textbf{\alph*.}]
  \item Démontrer qu'on a : $-\mathrm{e}^{-x}\le f(x)\le\mathrm{e}^{-x}$.
  \item Déterminer les coordonnées des points d'intersection de $(\mathcal{C}_f)$ avec les courbes
    d'équations $y=\mathrm{e}^{-x}$ et $y=-\mathrm{e}^{-x}$.
  \end{enumerate}
\item Sur $[0\,;\,2\pi]$, tracer dans le même repère les courbes d'équations $y=\mathrm{e}^{-x}$ et
  $y=-\mathrm{e}^{-x}$ puis la courbe $(\mathcal{C}_f)$.
\item Calculer l'aire de la partie du plan délimitée par $(\mathcal{C}_f)$ et la courbe d'équation
  $y=\mathrm{e}^{-x}$ sur $[0\,;\,2\pi]$. \textit{(On pourra utiliser la question 1.)}
\end{enumerate}

\subsection*{Partie B — Évaluation des compétences \hfill (15 points)}
\noindent\textbf{\underline{Situation :}}

\smallskip
\noindent Trois gisements de gaz $A$, $B$ et $C$ présentant chacun $100$ milliards de m\textsuperscript{3}
de quantité ont été découverts dans un pays. L'inauguration a eu lieu une certaine année
(année $0$) prise comme origine des temps $t$ (en années).

L'exploitation du gaz des gisements $A$ et $B$ avait commencé à la date $t=0$ et celle du
gisement $C$ légèrement avant. Seulement, à la date $t=1$, la quantité totale de gaz extraite
de chacun des gisements $A$ et $C$ était de $5{,}01$ milliards de m\textsuperscript{3}.

\begin{itemize}
\item Pour le gisement $A$ et à partir de la $2^{\text{e}}$ année, la quantité de gaz extraite
  chaque année augmente de $0{,}75$ milliard de m\textsuperscript{3} par rapport à celle de l'année
  précédente.
\item Pour les gisements $B$ et $C$, les ingénieurs pétrochimistes savent que si $q(t)$ est la
  quantité totale (en milliards de m\textsuperscript{3}) de gaz extraite de chacun de ces gisements à la
  date $t$, alors le taux d'extraction ou de consommation de gaz du gisement à cette date $t$
  est $q'(t)$ (milliards de m\textsuperscript{3} par an).
  \begin{itemize}
  \item Au niveau du gisement $B$, ce taux est $\left(\dfrac{1}{2t+1}+0{,}02t\right)$ milliard de
    m\textsuperscript{3} par an.
  \item Au niveau du gisement $C$, ces taux (aux dates $t$) sont proportionnels aux quantités de
    gaz extraites à ces dates. À la date $t=1$, ce taux était $5{,}01$ milliards de m\textsuperscript{3} par an.
  \end{itemize}
\end{itemize}

\noindent\textbf{\underline{Tâches :}}
\begin{enumerate}[label=\textbf{\arabic*.},leftmargin=2.2em]
\item En combien d'années le gisement $A$ s'épuisera-t-il ?
\item Combien d'années d'extraction suffiront à ce pays pour épuiser le contenu du gisement $B$ ?
\item Après l'inauguration, combien d'années faudra-t-il à ce pays pour vider le gisement $C$ de
  son contenu ?
\end{enumerate}

% ════════════════════════════════════════════════════
\corrige{de l'annale 2022}

\noindent\textbf{Partie A — Exercice 1.}

\smallskip
\noindent\textbf{1)} On résout $z^2+2-2\mathrm{i}\sqrt3=0$ :
\begin{align*}
z^2+2-2\mathrm{i}\sqrt3=0
&\iff z^2=-2+2\mathrm{i}\sqrt3\\
&\iff z^2=1+2\mathrm{i}\sqrt3-3\\
&\iff z^2=1+2\mathrm{i}\sqrt3+(\mathrm{i}\sqrt3)^2\\
&\iff z^2=(1+\mathrm{i}\sqrt3)^2\\
&\iff z=1+\mathrm{i}\sqrt3\quad\text{ou}\quad z=-(1+\mathrm{i}\sqrt3).
\end{align*}
L'ensemble des solutions est $S=\{\,1+\mathrm{i}\sqrt3\ ;\ -1-\mathrm{i}\sqrt3\,\}$.

\smallskip
\noindent\textbf{2.a)} Pour une similitude directe d'expression $z'=az+b$ avec $a\in\C^*\setminus\{1\}$
et $b\in\C$, le centre a pour affixe $\frac{b}{1-a}$, le rapport est $|a|$ et l'angle $\arg(a)$.

Ici $z'=(1-\mathrm{i}\sqrt3)z+0$ : comme $b=0$, le \textbf{centre de $s$ est $O$}.

Rapport : $|1-\mathrm{i}\sqrt3|=\sqrt{1^2+(\sqrt3)^2}=\sqrt4=2$.

Angle $\theta$ : $\begin{cases}\cos\theta=\dfrac12\\[2pt]\sin\theta=\dfrac{-\sqrt3}{2}\end{cases}$
donc $\theta=-\dfrac\pi3\ [2\pi]$.

Ainsi $s$ est la similitude directe de centre $O$, de rapport $2$ et d'angle $-\dfrac\pi3$.

\smallskip
\noindent\textbf{2.b)} Images de $A$ et $B$ par $s$.
\[ Z_{A'}=(1-\mathrm{i}\sqrt3)Z_A=(1-\mathrm{i}\sqrt3)(1+\mathrm{i}\sqrt3)=1+3=4=Z_F\ \Longrightarrow\ s(A)=F. \]
\[ Z_{B'}=(1-\mathrm{i}\sqrt3)Z_B=(1-\mathrm{i}\sqrt3)(-1-\mathrm{i}\sqrt3)=-(1+3)=-4=Z_G\ \Longrightarrow\ s(B)=G. \]

\smallskip
\noindent\textbf{3.a)} Le milieu de $[AB]$ est $O$ car $\dfrac{Z_A+Z_B}{2}=\dfrac{(1+\mathrm{i}\sqrt3)+(-1-\mathrm{i}\sqrt3)}{2}=0$.
Le centre de $(\mathcal{E})$ est donc $O$, et comme $s$ est une similitude de centre $O$, le centre
de $(\mathcal{E}')$ est aussi $O$. Une équation de $(\mathcal{E}')$ est donc de la forme
$\dfrac{x^2}{a^2}+\dfrac{y^2}{b^2}=1$.

Or l'ellipse $(\mathcal{E}')$ a pour foyers $F$ et $G$ (images de $A$ et $B$) :
\[ \begin{cases}c=OF=4\\[2pt]e=\dfrac ca=\dfrac12\end{cases}\Longrightarrow a=8\Longrightarrow a^2=64. \]
Puis $c^2=a^2-b^2\Rightarrow 16=64-b^2\Rightarrow b^2=48$ (et $b=\sqrt{48}=4\sqrt3$).

Une équation de $(\mathcal{E}')$ est donc : $\boxed{\dfrac{x^2}{64}+\dfrac{y^2}{48}=1.}$

\smallskip
\noindent\textbf{3.b)} On construit d'abord $(\mathcal{E}')$ (grand axe sur $(Ox)$, demi-axes $a=8$ et
$b=4\sqrt3\approx6{,}9$). La courbe $(\mathcal{E})$ s'obtient comme image de $(\mathcal{E}')$ par la similitude
$s^{-1}$ de centre $O$, de rapport $\frac12$ et d'angle $\frac\pi3$.

\begin{illus}
\begin{tikzpicture}[scale=0.42,>=Stealth]
  \draw[->,onbline] (-9,0)--(9,0) node[right]{$x$};
  \draw[->,onbline] (0,-8)--(0,8) node[above]{$y$};
  % ellipse E' : a=8, b=4sqrt3 ~ 6.93
  \draw[onbo,very thick] (0,0) ellipse (8 and 6.93);
  % ellipse E : image par s^{-1} (rapport 1/2, angle pi/3) -> demi-axes 4 et 2sqrt3, tournés de pi/3
  \draw[onbgreen,very thick,rotate=60] (0,0) ellipse (4 and 3.46);
  \node[onbo] at (6.3,5) {$(\mathcal{E}')$};
  \node[onbgreen] at (2.6,2.6) {$(\mathcal{E})$};
  \foreach \p/\xx/\yy/\pos in {F/4/0/below right, G/-4/0/below left}
    {\fill[onbblue] (\xx,\yy) circle (5pt); \node[onbblue] at (\xx,\yy) [\pos]{$\p$};}
  \fill[onbblue] (1,1.73) circle (4pt); \node[onbblue] at (1,1.73)[above right]{$A$};
  \fill[onbblue] (-1,-1.73) circle (4pt); \node[onbblue] at (-1,-1.73)[below left]{$B$};
  \fill (0,0) circle (4pt); \node at (0,0)[below left]{$O$};
\end{tikzpicture}
\end{illus}
\fig{$(\mathcal{E}')$ en orange (foyers $F$, $G$) et son image $(\mathcal{E})$ en vert (foyers $A$, $B$).}

\smallskip
\noindent\textbf{4)} Parmi les cinq points $O$, $A$, $B$, $F$, $G$, l'axe focal de $(\mathcal{E}')$ — qui est
l'axe $(Ox)$ — comprend les trois points $O$, $F$ et $G$.

Le tirage successif de deux points distincts parmi $5$ donne $A_5^2=5\times4=20$ choix
équiprobables. Le tirage de deux points distincts parmi les $3$ de l'axe focal donne
$A_3^2=3\times2=6$ choix favorables.

La probabilité cherchée est $\dfrac{6}{20}=\boxed{\dfrac{3}{10}}$.

\bigskip
\noindent\textbf{Exercice 2.}

\smallskip
\noindent\textbf{1.a)} On montre que $E_k$ est un sous-espace vectoriel de $E$.

$\bullet$ $f$ est un endomorphisme, donc $f(\vec0)=\vec0=k\vec0$, d'où $\vec0\in E_k$.

$\bullet$ Soient $\vec u,\vec v\in E_k$. Par linéarité de $f$ :
$f(\vec u+\vec v)=f(\vec u)+f(\vec v)=k\vec u+k\vec v=k(\vec u+\vec v)$, donc $\vec u+\vec v\in E_k$.

$\bullet$ Soient $\vec u\in E_k$ et $\lambda\in\R$ : $f(\lambda\vec u)=\lambda f(\vec u)=\lambda(k\vec u)=k(\lambda\vec u)$,
donc $\lambda\vec u\in E_k$.

Par conséquent, $E_k$ est un sous-espace vectoriel de $E$.

\smallskip
\noindent\textbf{1.b)} On suppose $f\circ f=2f$. Montrons que $\vec u\in\operatorname{Im}f\iff\vec u\in E_2$.

$\bullet$ \emph{Si $\vec u\in E_2$}, alors $f(\vec u)=2\vec u$, donc $\vec u=\frac12 f(\vec u)=f\!\left(\frac12\vec u\right)$,
ce qui montre que $\vec u\in\operatorname{Im}f$.

$\bullet$ \emph{Si $\vec u\in\operatorname{Im}f$}, il existe $\vec v\in E$ tel que $\vec u=f(\vec v)$. Alors
\[ f(\vec u)=f\big(f(\vec v)\big)=(f\circ f)(\vec v)=2f(\vec v)=2\vec u, \]
donc $\vec u\in E_2$.

Par conséquent, $\vec u\in\operatorname{Im}f$ si et seulement si $\vec u\in E_2$.

\smallskip
\noindent\textbf{2.a)} En additionnant membre à membre $(1)$ et $(2)$ :
\[ f(\vec\imath+\vec\jmath)+f(\vec\imath-\vec\jmath)=4\vec\imath,\qquad\text{or}\quad f(\vec\imath+\vec\jmath)+f(\vec\imath-\vec\jmath)=2f(\vec\imath), \]
d'où $2f(\vec\imath)=4\vec\imath$, soit $f(\vec\imath)=2\vec\imath$.

En soustrayant $(2)$ de $(1)$ : $f(\vec\imath+\vec\jmath)-f(\vec\imath-\vec\jmath)=4\vec\jmath=2f(\vec\jmath)$, d'où $f(\vec\jmath)=2\vec\jmath$.

Avec $(3)$ : $f(\vec\imath-\vec\jmath+\vec k)=\vec0\Rightarrow f(\vec\imath)-f(\vec\jmath)+f(\vec k)=\vec0\Rightarrow 2\vec\imath-2\vec\jmath+f(\vec k)=\vec0$,
d'où $f(\vec k)=-2\vec\imath+2\vec\jmath$.

\smallskip
\noindent\textbf{2.b)} Des relations $f(\vec\imath)=2\vec\imath$, $f(\vec\jmath)=2\vec\jmath$, $f(\vec k)=-2\vec\imath+2\vec\jmath$ :
\[ M=\begin{pmatrix}2&0&-2\\0&2&2\\0&0&0\end{pmatrix}. \]

\smallskip
\noindent\textbf{2.c)} Démontrer $f\circ f=2f$ revient à vérifier $M\times M=2M$ :
\[ M\times M=\begin{pmatrix}2&0&-2\\0&2&2\\0&0&0\end{pmatrix}\begin{pmatrix}2&0&-2\\0&2&2\\0&0&0\end{pmatrix}
=\begin{pmatrix}4&0&-4\\0&4&4\\0&0&0\end{pmatrix}=2\begin{pmatrix}2&0&-2\\0&2&2\\0&0&0\end{pmatrix}=2M. \]
Donc $M\times M=2M$, d'où $f\circ f=2f$.

\smallskip
\noindent\textbf{2.d)} Soit $\vec u(x\,;\,y\,;\,z)\in E$. Alors
\[ \vec u\in\operatorname{Ker}f\iff f(\vec u)=\vec0\iff\begin{cases}2x-2z=0\\2y+2z=0\end{cases}\iff\begin{cases}x=z\\y=-z.\end{cases} \]
Les coordonnées de $\vec u$ sont donc $(z\,;\,-z\,;\,z)$, soit $\vec u=z(\vec\imath-\vec\jmath+\vec k)$.
Par conséquent, $\operatorname{Ker}f$ est la droite vectorielle dont une base est $(\vec\imath-\vec\jmath+\vec k)$.

\smallskip
\noindent\textbf{2.e)} Par le théorème du rang, $\dim E=\dim(\operatorname{Im}f)+\dim(\operatorname{Ker}f)$,
soit $3=\dim(\operatorname{Im}f)+1$, donc $\dim(\operatorname{Im}f)=2$.

Comme $(\vec\imath,\vec\jmath,\vec k)$ est une base de $E$, la famille $\big(f(\vec\imath),f(\vec\jmath),f(\vec k)\big)$ engendre
$\operatorname{Im}f$. Or $f(\vec k)=-2\vec\imath+2\vec\jmath=-f(\vec\imath)+f(\vec\jmath)$, donc $\big(f(\vec\imath),f(\vec\jmath)\big)=(2\vec\imath,2\vec\jmath)$
suffit à engendrer $\operatorname{Im}f$.

Par conséquent, $\operatorname{Im}f$ est le plan vectoriel dont une base est $(\vec\imath,\vec\jmath)$.

\bigskip
\noindent\textbf{Exercice 3.}

\smallskip
\noindent\textbf{1)} Avec $f(x)=\mathrm{e}^{-x}\cos(x)$ :
\begin{align*}
f'(x)&=-\mathrm{e}^{-x}\cos(x)+\mathrm{e}^{-x}(-\sin(x))=-\mathrm{e}^{-x}\cos(x)-\mathrm{e}^{-x}\sin(x),\\
f''(x)&=-\big(\mathrm{e}^{-x}\cos(x)\big)'-\big(\mathrm{e}^{-x}\sin(x)\big)'\\
&=-\big(-\mathrm{e}^{-x}\cos(x)-\mathrm{e}^{-x}\sin(x)\big)-\big(-\mathrm{e}^{-x}\sin(x)+\mathrm{e}^{-x}\cos(x)\big)\\
&=\mathrm{e}^{-x}\cos(x)+\mathrm{e}^{-x}\sin(x)+\mathrm{e}^{-x}\sin(x)-\mathrm{e}^{-x}\cos(x)=2\mathrm{e}^{-x}\sin(x).
\end{align*}
Dès lors :
\begin{align*}
f''(x)+2f'(x)+2f(x)&=2\mathrm{e}^{-x}\sin(x)+2\big(-\mathrm{e}^{-x}\cos(x)-\mathrm{e}^{-x}\sin(x)\big)+2\mathrm{e}^{-x}\cos(x)\\
&=2\mathrm{e}^{-x}\sin(x)-2\mathrm{e}^{-x}\cos(x)-2\mathrm{e}^{-x}\sin(x)+2\mathrm{e}^{-x}\cos(x)=0.
\end{align*}
Donc $f''(x)+2f'(x)+2f(x)=0$.

\smallskip
\noindent\textbf{2)} On a $f'(x)=\mathrm{e}^{-x}\big(-\cos(x)-\sin(x)\big)$. Comme $\mathrm{e}^{-x}>0$, le signe de
$f'(x)$ est celui de $-\cos(x)-\sin(x)$. On résout sur $[0\,;\,2\pi]$ :
\[ -\cos(x)-\sin(x)=0\iff\sin(x)=-\cos(x)\iff\tan(x)=-1\iff x=\frac{3\pi}{4}\ \text{ou}\ x=\frac{7\pi}{4}. \]
Par continuité, on étudie le signe sur chaque intervalle :
\[ x=0:\ -1<0\ ;\qquad x=\pi:\ -(-1)-0=1>0\ ;\qquad x=2\pi:\ -1<0. \]
Calculs préliminaires :
\[ f(0)=1,\quad f\!\left(\tfrac{3\pi}{4}\right)=-\tfrac{\sqrt2}{2}\mathrm{e}^{-3\pi/4},\quad
f\!\left(\tfrac{7\pi}{4}\right)=\tfrac{\sqrt2}{2}\mathrm{e}^{-7\pi/4},\quad f(2\pi)=\mathrm{e}^{-2\pi}. \]

\begin{center}\renewcommand{\arraystretch}{1.6}
\begin{tabular}{|c|ccccccc|}
\hline
$x$ & $0$ && $\dfrac{3\pi}{4}$ && $\dfrac{7\pi}{4}$ && $2\pi$\\
\hline
$f'(x)$ & & $-$ & $0$ & $+$ & $0$ & $-$ &\\
\hline
$f(x)$ & $1$ & $\searrow$ & $-\dfrac{\sqrt2}{2}\mathrm{e}^{-3\pi/4}$ & $\nearrow$ & $\dfrac{\sqrt2}{2}\mathrm{e}^{-7\pi/4}$ & $\searrow$ & $\mathrm{e}^{-2\pi}$\\
\hline
\end{tabular}
\end{center}

\smallskip
\noindent\textbf{3.a)} Pour tout $x\in[0\,;\,2\pi]$ : $\mathrm{e}^{-x}>0$ et $-1\le\cos(x)\le1$, d'où en
multipliant par $\mathrm{e}^{-x}>0$ : $-\mathrm{e}^{-x}\le\mathrm{e}^{-x}\cos(x)\le\mathrm{e}^{-x}$, soit $-\mathrm{e}^{-x}\le f(x)\le\mathrm{e}^{-x}$.

\smallskip
\noindent\textbf{3.b)} \emph{Avec $y=\mathrm{e}^{-x}$ :}
\[ f(x)=\mathrm{e}^{-x}\iff\mathrm{e}^{-x}\cos(x)=\mathrm{e}^{-x}\iff\cos(x)=1\iff x=0\ \text{ou}\ x=2\pi. \]
Les points d'intersection sont $(0\,;\,1)$ et $(2\pi\,;\,\mathrm{e}^{-2\pi})$.

\emph{Avec $y=-\mathrm{e}^{-x}$ :}
\[ f(x)=-\mathrm{e}^{-x}\iff\cos(x)=-1\iff x=\pi. \]
Le point d'intersection est $(\pi\,;\,-\mathrm{e}^{-\pi})$.

\smallskip
\noindent\textbf{4)} Tracé : $y=\mathrm{e}^{-x}$ (vert), $y=-\mathrm{e}^{-x}$ (bleu), $(\mathcal{C}_f)$ (rouge),
oscillant entre les deux enveloppes et les touchant aux points trouvés.

\begin{illus}
\begin{tikzpicture}[scale=1,>=Stealth,domain=0:6.4,samples=120]
  \draw[->,onbline] (-0.2,0)--(6.7,0) node[right]{$x$};
  \draw[->,onbline] (0,-1.05)--(0,1.15) node[above]{$y$};
  \draw[dashed,onbline] (3.1416,0.04)--(3.1416,-0.043);
  \draw[onbgreen,thick] plot (\x,{exp(-\x)}) node[right]{\footnotesize$y=\mathrm{e}^{-x}$};
  \draw[onbblue,thick] plot (\x,{-exp(-\x)}) node[right]{\footnotesize$y=-\mathrm{e}^{-x}$};
  \draw[onbo,very thick] plot (\x,{exp(-\x)*cos(\x r)});
  \node[onbo] at (1.2,0.55) {$(\mathcal{C}_f)$};
\end{tikzpicture}
\end{illus}
\fig{Encadrement de $(\mathcal{C}_f)$ entre $y=\mathrm{e}^{-x}$ et $y=-\mathrm{e}^{-x}$ sur $[0\,;\,2\pi]$.}

\smallskip
\noindent\textbf{5)} Sur $[0\,;\,2\pi]$, $\mathrm{e}^{-x}\ge f(x)$, donc l'aire (en unités d'aire) est
$\displaystyle\int_0^{2\pi}\big(\mathrm{e}^{-x}-f(x)\big)\dd x=\int_0^{2\pi}\mathrm{e}^{-x}\dd x-\int_0^{2\pi}f(x)\dd x$.

$\bullet$ $\displaystyle\int_0^{2\pi}\mathrm{e}^{-x}\dd x=\big[-\mathrm{e}^{-x}\big]_0^{2\pi}=1-\mathrm{e}^{-2\pi}.$

$\bullet$ D'après la question 1, $2f(x)=-f''(x)-2f'(x)$, soit $f(x)=-\tfrac12 f''(x)-f'(x)$. Donc
\begin{align*}
\int_0^{2\pi}f(x)\dd x&=-\tfrac12\big[f'(x)\big]_0^{2\pi}-\big[f(x)\big]_0^{2\pi}
=\tfrac12\Big[\mathrm{e}^{-x}\big(\sin(x)-\cos(x)\big)\Big]_0^{2\pi}\\
&=\tfrac12\Big[\mathrm{e}^{-2\pi}(0-1)-1\cdot(0-1)\Big]=\tfrac12\big(1-\mathrm{e}^{-2\pi}\big).
\end{align*}
D'où
\[ \int_0^{2\pi}\big(\mathrm{e}^{-x}-f(x)\big)\dd x=(1-\mathrm{e}^{-2\pi})-\tfrac12(1-\mathrm{e}^{-2\pi})=\tfrac12\big(1-\mathrm{e}^{-2\pi}\big)\ \text{u.a.} \]
L'unité d'aire vaut $2\times4=8$~cm\textsuperscript{2}. L'aire cherchée est donc
\[ 8\times\tfrac12\big(1-\mathrm{e}^{-2\pi}\big)=\boxed{4\big(1-\mathrm{e}^{-2\pi}\big)\ \text{cm}^2.} \]

\bigskip
\noindent\textbf{Partie B — Évaluation des compétences.}

\smallskip
\noindent\textbf{Tâche 1 (gisement $A$).} Soit $u_n$ la quantité de gaz extraite l'année $n$ (en
milliards de m\textsuperscript{3}). On a $u_1=5{,}01$ et $u_{n+1}=u_n+0{,}75$ : $(u_n)$ est arithmétique de
raison $r=0{,}75$ et de premier terme $u_1=5{,}01$.
\[ u_n=u_1+(n-1)r=5{,}01+0{,}75(n-1)=4{,}26+0{,}75n. \]
La quantité totale extraite au bout de $n$ années est
\[ S_n=n\cdot\frac{u_1+u_n}{2}=n\cdot\frac{5{,}01+4{,}26+0{,}75n}{2}=n(4{,}635+0{,}375n)=0{,}375n^2+4{,}635n. \]
Le gisement s'épuise dès que $S_n\ge100$, soit $0{,}375n^2+4{,}635n-100\ge0$.
\[ \Delta=4{,}635^2-4(0{,}375)(-100)=21{,}483225+150=171{,}483225>0, \]
\[ n_1=\frac{-4{,}635-\sqrt{171{,}483225}}{0{,}75}\approx-23{,}64,\qquad
n_2=\frac{-4{,}635+\sqrt{171{,}483225}}{0{,}75}\approx11{,}28. \]
Le trinôme est positif pour $n\ge n_2\approx11{,}28$. Le plus petit entier naturel convenable est
$n=12$. \textbf{Le gisement $A$ s'épuisera en $12$ ans.}

\smallskip
\noindent\textbf{Tâche 2 (gisement $B$).} On a $q'(t)=\dfrac{1}{2t+1}+0{,}02t$, d'où
\[ q(t)=\tfrac12\ln(2t+1)+0{,}01t^2+k\quad(k\in\R). \]
Comme $q(0)=0$ : $\tfrac12\ln(1)+0+k=0\Rightarrow k=0$. Donc $q(t)=\tfrac12\ln(2t+1)+0{,}01t^2$.

$q$ est continue et strictement croissante sur $]0\,;\,+\infty[$ (somme d'une fonction croissante
et de $0{,}01t^2$ croissante ; $q'(t)>0$). Or
\[ q(98)=\tfrac12\ln(197)+0{,}01(98)^2\approx98{,}68<100,\qquad q(99)=\tfrac12\ln(199)+0{,}01(99)^2\approx100{,}66>100. \]
D'après le corollaire du théorème des valeurs intermédiaires, l'équation $q(t)=100$ admet une
unique solution dans $]98\,;\,99[$. \textbf{Le gisement $B$ sera épuisé au cours de la $99^{\text{e}}$ année.}

\smallskip
\noindent\textbf{Tâche 3 (gisement $C$).} Les taux étant proportionnels aux quantités extraites,
$\dfrac{q'(t)}{q(t)}$ est constant ; à $t=1$, $\dfrac{q'(1)}{q(1)}=\dfrac{5{,}01}{5{,}01}=1$, donc $q'(t)=q(t)$.
Les solutions sont $q(t)=C\mathrm{e}^{t}$ ($C$ constante). Avec $q(1)=5{,}01$ :
\[ C\mathrm{e}=5{,}01\Rightarrow C=5{,}01\,\mathrm{e}^{-1},\qquad\text{d'où}\quad q(t)=5{,}01\,\mathrm{e}^{t-1}. \]
Le gisement est vidé quand $q(t)=100$ :
\[ 5{,}01\,\mathrm{e}^{t-1}=100\iff\mathrm{e}^{t-1}=\frac{100}{5{,}01}\iff t-1=\ln\!\left(\frac{100}{5{,}01}\right)\iff
t=1+\ln\!\left(\frac{100}{5{,}01}\right)\approx3{,}99. \]
\textbf{Après l'inauguration, il faudra environ $4$ ans pour vider le gisement $C$.}
